\scanpage{81}

% upper, lines 1–8
enfin trois \deleted{familles} \inserted{classes de conjugaison} de ss-gpes \deleted{conjugués} d'indice $6$, \inserted{i.e. d'ordre $2$} comprenant un gpe invariant (i.e. formant une seule classe), savoir le centre, qui est le stabilisateur des él. de $H_{r'}$, et $2$ \deleted{familles} \inserted{classes} de \inserted{$3$ ss-}groupes chacunes \deleted{\unsure{comprenant}} resp. stabilisateurs des él. de $H_s$ et $H_{s*}$ : les \deleted{\unsure{$\varepsilon$}} $u$ \inserted{et les $zu$} ($z$ él. non triv. du centre, i.e. gén. du facteur $\mathbb Z/2\mathbb Z=\mathfrak S_2$ de $\mathfrak S_2\times\mathfrak S_3$, $u$ el. d'ordre $2$ de $\mathfrak S_3$) \deleted{et} \deleted{\unsure{les $zu$}}

% middle, lines 9–15
Les classes de conj. de ss-gpes de $D_6$ correspondent encore \inserted{$1-1$} aux \deleted{classes} espaces hom. du diagramme marqué, et on constate que les autom. de ces espaces homogènes sont ceux marqués sur le diagramme (une flèche $\xrightarrow{2}$ i.e. d'ordre $3$ impliquant qu'il y a à la source un autom. inv. (ss pt fixe) --).

% archive sheet number: 80
